\section*{Capítulo \ref{ch:b3:field-lab-methods-3} --- Métodos de campo y de laboratorio III}
% Values from the Book 5 scratch calculator calc/ch31.py.

\begin{solution}{exo:b3:field-lab-methods-3:1}
(a)~$\qty{0.01}{\milli\gal} = \qty{e-7}{m.s^{-2}}$, y
$\num{e-7}/\num{9.81} = \num{1.0e-8}$: una parte en cien millones.

(b)~$0.01/0.3086 = \qty{0.032}{m}$, es decir, unos \qty{3.2}{cm}.

(c)~Juntos, los dos efectos cambian la gravedad en
\qty{0.1967}{\milli\gal.m^{-1}}, la diferencia $0.3086 - 0.1119$, de modo que
la altura debe cambiar en $0.01/0.1967 = \qty{5.1}{cm}$.
\end{solution}

\begin{solution}{exo:b3:field-lab-methods-3:2}
$\Delta m = \tfrac43\pi(\qty{50}{m})^3\times\qty{500}{kg.m^{-3}} =
\qty{2.62e8}{kg}$. El pico es $G\Delta m/z^2 =
\num{6.674e-11}\times\num{2.62e8}/\num{e4} = \qty{1.75e-6}{m.s^{-2}}$,
es decir, \qty{0.175}{\milli\gal}. La semianchura es $x_{1/2} = z/1.305 =
\qty{77}{m}$.
\end{solution}

\begin{solution}{exo:b3:field-lab-methods-3:3}
(a)~Con $AM = MN = NB = a$, $K = 2\pi/(1/a - 1/2a - 1/2a + 1/a) =
2\pi a = \qty{62.83}{m}$, y $\rho_a = 62.83\times0.239/0.1 =
\qty{150}{\ohm\metre}$. (b)~$K = \pi(15^2 - 5^2)/(2\times5) = \pi\times
200/10 = \qty{62.83}{m}$. Los electrodos quedan entonces en $0$, $10$, $20$
y \qty{30}{m}, en el orden AMNB: se trata de la configuración Wenner de separación
\qty{10}{m}.
\end{solution}

\begin{solution}{exo:b3:field-lab-methods-3:4}
(a)~$V = 0.2998/\sqrt5 = \qty{0.134}{m.ns^{-1}}$. (b)~$t =
2\times10/0.134 = \qty{149}{ns}$. (c)~$\delta = 0.0053\sqrt5/\num{e-3} =
\qty{11.9}{m}$, y una profundidad útil de unos \qty{36}{m}. (d)~En la arcilla
$\delta = 0.0053\sqrt{20}/0.01 = \qty{2.4}{m}$, una profundidad útil de unos
\qty{7}{m}, y una velocidad de solo \qty{0.067}{m.ns^{-1}}. El
aproximación sigue siendo válida allí: a \qty{100}{MHz},
$\omega\varepsilon = 2\pi\times\num{e8}\times20\times\num{8.85e-12} =
\qty{0.11}{S.m^{-1}}$, diez veces $\sigma$.
\end{solution}

\begin{solution}{exo:b3:field-lab-methods-3:5}
(a)~El cilindro tiene $\lambda = \pi R^2\Delta\rho$ por unidad de longitud, y
fuera de él su atracción es la de una masa lineal sobre su eje (teorema de Gauss
sobre un cilindro coaxial): $\Delta g = 2G\pi R^2\Delta\rho\,z/(x^2
+ z^2)$. La mitad del pico se alcanza cuando $x^2 + z^2 = 2z^2$, por tanto, $z =
x_{1/2}$. (b)~La esfera está a $1.305\times100 = \qty{130}{m}$, el
cilindro a \qty{100}{m}: el mismo perfil, interpretado como esfera, sitúa el
cuerpo \qty{30}{\percent} más profundo. La anomalía de una esfera decae más deprisa con
la distancia, así que una esfera debe estar más profunda para ser igual de ancha.
\end{solution}

\begin{solution}{exo:b3:field-lab-methods-3:6}
(a)~Cada punto representa $\qty{100}{m}\times\qty{100}{m} =
\qty{e4}{m^2}$, así que $\iint\Delta g\,\mathrm{d}A \approx
\num{-438e-5}\times\num{e4} = \qty{-43.8}{m^3.s^{-2}}$, y $\Delta m
= -43.8/(2\pi\times\num{6.674e-11}) = \qty{-1.04e11}{kg}$. (b)~$1.04/1.74
= \qty{60}{\percent}$. El resto está en los flancos de la anomalía fuera
del mapa: el granito alcanza una profundidad de más de un kilómetro, y la
anomalía de su parte profunda decae lentamente con la distancia. Ampliar el mapa
lo suficiente, o usar un modelo, permite recuperarla.
\end{solution}

\begin{solution}{exo:b3:field-lab-methods-3:7}
(a)~Para un dique delgado la anomalía cae a la mitad de su pico en $x = \pm z$:
su anchura total a mitad del pico es $2z = \qty{60}{m}$, y tres
lecturas a través de él requieren una separación de como máximo \qty{20}{m}. (b)~$2\times
220/3 = \qty{147}{m}$. (c)~La separación entre estaciones necesaria escala con la
profundidad del cuerpo de interés más somero: una prospección diseñada para el
dique profundo, a \qty{150}{m}, submuestrearía la anomalía del somero, que
tiene solo \qty{60}{m} de anchura; una diseñada para el cuerpo somero también detecta el
profundo, con muchas más lecturas de las necesarias.
\end{solution}

\begin{solution}{exo:b3:field-lab-methods-3:8}
(a)~$F = 80/8 = 10$ y $\phi = F^{-1/m} = 10^{-1/2} = \num{0.316}$.
(b)~$10^{-1/1.8} = \num{0.278}$ y $10^{-1/2.2} = \num{0.351}$. Un cambio
de \num{0.2} en $m$ desplaza la porosidad unos cuatro puntos porcentuales;
las resistividades, medidas con una precisión de unos pocos por ciento, la desplazan mucho menos. El
exponente de cementación controla la incertidumbre, y debe calibrarse
con muestras de la misma roca.
\end{solution}

\begin{solution}{exo:b3:field-lab-methods-3:9}
(a)~Para $s = 2h_1$, $(2nh_1/s)^2 = n^2$, así que los términos son $k^n(1 +
n^2)^{-3/2}$: $0.3536k$, $0.0894k^2$, $0.0316k^3$. Para $k = 0.5$:
$\num{0.1768}$, $\num{0.0224}$, $\num{0.0040}$, y $\rho_a/\rho_1 \approx
1 + 2\times0.2031 = \num{1.406}$ (la suma completa da \num{1.409}). Para $k
= -0.5$ los términos alternan de signo, $-0.1768 + 0.0224 - 0.0040$, y
$\rho_a/\rho_1 \approx \num{0.683}$ (suma completa \num{0.685}). (b)~Para $s \ll
h_1$ todos los términos se anulan y $\rho_a \to \rho_1$; para $s \gg h_1$ cada
corchete tiende a~1 y $1 + 2\sum k^n = 1 + 2k/(1 - k) = (1 + k)/(1 -
k) = \rho_2/\rho_1$. (c)~$\rho_2/\rho_1 = (1 + k)/(1 - k) = 3$ y $1/3$.
\end{solution}

\begin{solution}{exo:b3:field-lab-methods-3:10}
(a)~En el orden AMNB con separación $a$: $AM = a$, $BM = 2a$, $AN = 2a$,
$BN = a$, así que el corchete vale $2/a - 1/a = 1/a$ y $K = 2\pi a$.
(b)~Con A y B en $\mp s$ y M y N en $\mp l$: $AM = BN = s - l$,
$BM = AN = s + l$, el corchete vale $2/(s - l) - 2/(s + l) = 4l/(s^2 -
l^2)$, y $K = \pi(s^2 - l^2)/(2l)$. Cuando $l \ll s$, $\Delta V_{MN}
\approx 2l\,E$ con $E = 2(-\mathrm{d}V/\mathrm{d}r)$ en el centro, el
campo de ambos electrodos de corriente; entonces $\rho_a = K\Delta V/I \to
\pi s^2/(2l)\times 4l(-\mathrm{d}V/\mathrm{d}r)/I = 2\pi s^2(-
\mathrm{d}V/\mathrm{d}r)/I$. (c)~La razón entre el factor aproximado y el
factor exacto es $s^2/(s^2 - l^2) = 1/(1 - 0.01) = 1.010$: el 1 por ciento.
\end{solution}

\begin{solution}{exo:b3:field-lab-methods-3:11}
(a)~Lectura: \qty{0.010}{\milli\gal}. Altura:
$0.1967\times0.05 = \qty{0.010}{\milli\gal}$. Densidad: la corrección de Bouguer
es $0.04191\,h\,d$ con $d$ en \unit{g.cm^{-3}}, así que
$0.04191\times50\times0.1 = \qty{0.210}{\milli\gal}$. Topografía:
\qty{0.200}{\milli\gal}. Total: $\sqrt{0.010^2 + 0.010^2 + 0.210^2 +
0.200^2} = \qty{0.29}{\milli\gal}$. (b)~La densidad y la topografía
dominan, cada una unas veinte veces el error de lectura. Para detectar
\qty{0.1}{\milli\gal}, la densidad de la losa debe medirse en
muestras o ajustarse a los datos (el valor que menos refleja la
topografía), y la \omterm{def:b3:field-lab-methods-3:terrain}{corrección topográfica} debe calcularse a partir de un modelo digital de elevaciones detallado.
(c)~$\sqrt{0.010^2 + 0.010^2 + 0.210^2} = \qty{0.21}{\milli\gal}$.
\end{solution}

\begin{solution}{exo:b3:field-lab-methods-3:12}
(a)~La masa que ajusta la lectura central es $\Delta m = \Delta g(0)z^2/G
= \num{-4.16e-5}z^2/\num{6.674e-11}$: \qty{-5.61e10}{kg} a
\qty{300}{m}, \qty{-9.97e10}{kg} a \qty{400}{m}, \qty{-1.56e11}{kg} a
\qty{500}{m}. Las lecturas predichas en $x = \pm300$ y $\pm600$ son
\num{-1.47} y \num{-0.37}, \num{-2.13} y \num{-0.71}, \num{-2.62}
y \qty{-1.09}{\milli\gal}. Los \omterm{def:b3:field-lab-methods-3:inversion-practice}{desajustes} son $\chi^2 = 1238$, $3.7$ y
$856$. (b)~Solo el modelo a \qty{400}{m}, con $\chi^2 = 3.7$ para cinco
datos, se ajusta dentro de los errores. (c)~Sí: cualquier cuerpo de simetría esférica
con la misma masa y el mismo centro produce exactamente la misma anomalía, y otras
formas pueden quedar dentro de los errores. Una densidad medida en muestras, en
una perforación u otro método lo decidiría.
\end{solution}

\begin{solution}{pb:b3:field-lab-methods-3:1}
\textbf{1.} $(2846.708 - 2846.512)/\qty{2.0}{h} = \qty{0.098}{\milli\gal.h^{-1}}$.

\textbf{2.} $2851.075 - 0.098\times0.5 - 2846.512 = \qty{4.514}{\milli\gal}$;
$2850.497 - 0.088 - 2846.512 = \qty{3.897}{\milli\gal}$;
$2849.392 - 0.127 - 2846.512 = \qty{2.753}{\milli\gal}$.

\textbf{3.} Las alturas sobre la base son $-27.1$, $-30.6$ y
\qty{-18.1}{m}; sumando $0.1967\,\Delta h$ se obtienen las anomalías de Bouguer
$\num{-0.817}$, $\num{-2.122}$ y \qty{-0.808}{\milli\gal}. Las
estaciones están más bajas que la base y más cerca del centro de la Tierra,
por lo que sus lecturas brutas son más altas; la corrección elimina eso y
deja un mínimo sobre el granito.

\textbf{4.} El exceso $0.098 - 0.040 = \qty{0.058}{\milli\gal.h^{-1}}$,
\qty{0.12}{\milli\gal} en el circuito, es la marea terrestre. Durante dos
horas varía casi linealmente, así que una recta que pase por las dos
lecturas en la base la elimina junto con la deriva.

\textbf{5.} Los errores son $\num{0.008}$, $\num{0.009}$ y
\qty{0.017}{\milli\gal}. El redondeo de cada altura a \qty{0.1}{m} puede hacer
que una diferencia de alturas sea errónea en hasta \qty{0.1}{m}, es decir,
\qty{0.020}{\milli\gal}: el redondeo explica los errores, y la marea
que queda tras la recta es inferior a \qty{0.001}{\milli\gal}.

\textbf{6.} $0.01/0.1967 = \qty{5.1}{cm}$. Los \qty{3.3}{cm} de los manuales corresponden a
$0.01/0.3086$: la anomalía de aire libre no depende de la densidad
elegida para la corrección de Bouguer, y la altura debe servir también para ella.

\textbf{7.} $x_{1/2} = (3.952 - 3.048)/2 = \qty{0.452}{km}$, y $z = 1.305
\times452 = \qty{590}{m}$ bajo el terreno.

\textbf{8.} $\Delta m = \Delta g\,z^2/G = \num{-2.293e-5}\times590^2/
\num{6.674e-11} = \qty{-1.20e11}{kg}$. Con $\Delta\rho =
\qty{-130}{kg.m^{-3}}$, su volumen es \qty{9.2e8}{m^3} y su radio
\qty{603}{m}, mayor que la profundidad de su centro: la esfera sobresaldría
del terreno. Ninguna esfera con el contraste de densidad del granito
encaja; el cuerpo debe llegar cerca de la superficie.

\textbf{9.} El granito se extiende desde la superficie hasta una gran profundidad. Su parte somera
, cerca de las estaciones, hace estrecho el máximo, y la regla, que
solo conoce la anchura, coloca una única esfera a una profundidad moderada; la
parte profunda añade flancos anchos que la esfera no tiene (la figura del
perfil), y la masa que contienen queda sin contabilizar.

\textbf{10.} $2\pi G\Delta\rho = 2\pi\times\num{6.674e-11}\times(-130) =
\qty{-5.45e-8}{s^{-2}}$. Con $z_1 = 0$, $z_2 = \qty{1645}{m}$, $a =
\qty{461}{m}$: el corchete vale $1645 + 461 - 1708.4 = \qty{397.6}{m}$,
y $\Delta g = \qty{-2.17}{\milli\gal}$. Para una profundidad ilimitada el corchete
tiende a $a$, y $\Delta g = \qty{-2.51}{\milli\gal}$.

\textbf{11.} El corchete debe valer $2.29/\num{5.45e-3} = \qty{420}{m}$, lo que
da un fondo a unos \qty{2.6}{km}. Cualquier fondo entre unos
\qty{1.6}{km} y una profundidad ilimitada cambia el máximo en unas pocas décimas
de miligal solamente, y un cuerpo algo más ancho o más denso da el mismo
máximo con un fondo más somero: el fondo no queda resuelto.

\textbf{12.} Con $z_1 = \qty{6}{m}$: $1639 + 461.0 - 1708.4 = 391.6$, y
$\Delta g = \qty{-2.14}{\milli\gal}$, \qty{0.03}{\milli\gal} menos que
con el techo en el terreno. Esto queda por debajo de los errores del terreno y de la densidad
de la campaña: la gravedad no puede distinguir si el granito alcanza
la superficie bajo los aluviones. El sondeo sí puede.

\textbf{13.} $K = \pi(20^2 - 1^2)/(2\times1) = \qty{626.7}{m}$, y $\Delta V =
\rho_a I/K = 821\times0.050/626.7 = \qty{65.5}{mV}$.

\textbf{14.} $V = 2\times1/\sqrt{17.1^2 - 10.7^2} = 2/13.34 =
\qty{0.150}{m.ns^{-1}}$, $\varepsilon_r = (0.2998/0.150)^2 = 4.0$, una
arena seca; la tubería está a $0.150\times10.7/2 = \qty{0.80}{m}$ de profundidad.

\textbf{15.} $0.150\times20.0/2 = \qty{1.50}{m}$.

\textbf{16.} $20/821 = \num{0.0244}$, $30/1309 = \num{0.0229}$, $40/1684 =
\num{0.0238}$: $S \approx \qty{0.0237}{S}$.

\textbf{17.} La capa seca aporta $1.5/600 = \qty{0.0025}{S}$, la
saturada $\qty{0.0212}{S}$, así que su espesor es $0.0212\times222 =
\qty{4.7}{m}$, y el granito yace $1.5 + 4.7 = \qty{6.2}{m}$ por debajo de
la superficie. La estimación es alta: el granito no es infinitamente
resistivo, así que a estas separaciones $\rho_a$ aún queda por debajo de la recta
$s/S$, y $s/\rho_a$ sobrestima $S$. (El modelo de la región tiene
\qty{6}{m} de aluvión.)

\textbf{18.} $\phi = (20/222)^{1/2} = \num{0.30}$. Con $m = 1.3$, $\phi =
0.0901^{1/1.3} = \num{0.16}$; con $m = 2.8$, $\phi = \num{0.42}$.

\textbf{19.} Al restar la base de la línea se obtiene, de oeste a este:
3.4, 4.1, 4.8, 5.9, 7.4, 9.2, 11.5, 14.2, 17.2, 19.8, 21.0, 20.1, 17.6,
14.5, 11.5, 9.1, 7.2, 5.7, 4.6, 3.8 y \qty{3.1}{nT}.

\textbf{20.} Sin corregir, el máximo es \qty{28.0}{nT} en \qty{2.20}{km}, y
la mitad, \qty{14.0}{nT}, nunca se alcanza hacia el este: la línea termina
en \qty{15.2}{nT}, porque el campo subió \qty{12}{nT} durante el
recorrido.

\textbf{21.} El máximo es \qty{21.0}{nT} en \qty{2.20}{km}; la mitad es
\qty{10.5}{nT}. Oeste: entre \num{9.2} en \num{1.95} y \num{11.5} en
\qty{2.00}{km}, en $1.95 + 0.05\times1.3/2.3 = \qty{1.978}{km}$. Este:
entre \num{11.5} en \num{2.40} y \num{9.1} en \qty{2.45}{km}, en
$2.40 + 0.05\times1.0/2.4 = \qty{2.421}{km}$. La semianchura es $(2.421
- 1.978)/2 = \qty{0.221}{km}$.

\textbf{22.} Una anomalía simétrica en un campo inclinado \ang{65} significa que solo
actúa la magnetización vertical: el cuerpo tiene rumbo según el
meridiano magnético, norte--sur. La regla de
la \cref{prop:b3:field-lab-methods-3:dyke} supone un dique delgado, vertical y
bidimensional que se extiende a gran profundidad, magnetizado por inducción, cuya
parte superior actúa como una línea de carga magnética.

\textbf{23.} $\kappa t = 2\pi z\,\Delta T_\mathrm{max}/(F\sin^2 I) = 2\pi\times
221\times21.0/(48000\times0.8214) = \qty{0.74}{m}$. Con $\kappa =
0.03$, $t = \qty{25}{m}$. Solo el producto está fijado: un dique con el doble de
espesor y la mitad de susceptibilidad da la misma anomalía.

\textbf{24.} La anomalía tiene \qty{443}{m} de anchura a mitad de su máximo y las estaciones
están separadas \qty{50}{m}: nueve lecturas sobre ella, tres veces el
mínimo de la regla.

\textbf{25.} La parte superior del dique está
\textbf{\qty[text-series-to-math=true]{220}{m}} bajo el terreno. El
modelo geológico sitúa la discordancia a \qty{222}{m} de profundidad en este punto:
el dique asciende hasta la base de la cobertera y se corta allí, más antiguo
que la arenisca triásica que lo sepulta.
\end{solution}
